\section{The end layers and the global mode}\label{sec:ends}

Near the plus end, $A_N=N-j$ with $j$ small means that the run in progress at time $N$
is a plus run of length of order $N$. Its probability is of order $T_\plus(N)$, and the
only question is where it started. It starts right after the $j$-th minus step exactly
when a minus run ends there, which gives the renewal mass function $u_\minus(j)$. The
next theorem makes this precise at fixed $j$ under regular variation alone, and
uniformly up to distance $(1-\delta)N$ from the end under Doney's condition.

\begin{introtheorem}[the end layers]\label{thm:ends}
Assume (BAL) and (SPAN).
\begin{enumerate}[(a)]
\item For each fixed $j\ge0$, a fresh plus start has
$P_N(N-j)=T_\plus(N)u_\minus(j)+o(T_\plus(N))$, and a fresh minus start has
$P_N(N-j)=T_\plus(N)u_\minus(j)\ones_{\{j\ge1\}}+o(T_\plus(N))$.
\item Assume moreover (LB$_\plus$) and (SRT$_\minus$). Then for a fresh plus start and
every $\delta\in(0,1)$, uniformly in $0\le j\le(1-\delta)N$,
\[
 P_N(N-j)=T_\plus(N)\bigl[u_\minus(j)\,\Phi_b(j/N)+o(\omega_\minus(j))\bigr],
\]
with $\Phi_b$ as in~\eqref{eq:Phi}. For a fresh minus start the same holds with
$u_\minus(j)\ones_{\{j\ge1\}}$ in place of $u_\minus(j)$.
\item Assume (LB$_\pm$) and $T_\pm(n)\sim c_\pm n^{-\alpha}$ with $c_\pm>0$. Then for a
fresh start of either sign and every sequence $J_N\to\infty$ with $J_N\le N/2$,
\[
 P_N(k)=N^{-1}\La(k/N)\,(1+o(1))\qquad\text{uniformly over }J_N\le k\le N-J_N .
\]
\end{enumerate}
The same statements hold at the minus end with the signs exchanged.
\end{introtheorem}

Part (a) has a simple reading. The event $A_N=N-j$ with $j$ fixed requires a plus run of length
of order $N$. On this event the first $N$ steps contain at most $2j+2$ runs, since the $j$
minus steps lie in at most $j+1$ minus runs, and the proof shows that the long run is, with
probability close to $1$, the run in progress at time $N$. The tail of
that run gives the factor $T_\plus(N)$. Before it, the minus steps that have been taken are
exactly $j$, and they are the union of complete minus runs. So $j$ must be a renewal epoch of
the minus runs, which has probability $u_\minus(j)$. For a minus start the epoch $j=0$ is
excluded, since the walk begins with a minus run. In the continuous-time model with one law
and a fair start, Wang, Schulz, Deng and Barkai~\cite[Section~VII.A]{wsdb2018} found an atom
with a boundary layer for the occupation time, and (a) is a lattice version of this layer for
two laws at a fixed time.

In (b) the factor $\Phi_b$ is what makes the layer and the bulk agree. Under
(SRT$_\minus$), uniformly for $\varepsilon N\le j\le(1-\delta)N$,
\begin{equation}\label{eq:layer-bulk}
 T_\plus(N)\,u_\minus(j)\,\Phi_b(j/N)=N^{-1}\La(1-j/N)\,(1+o(1)).
\end{equation}
Indeed, with $y=j/N$, (SRT$_\minus$) and (R1) give
$T_\plus(N)u_\minus(j)=\frac{\sin\pi\alpha}{\pi b}\,y^{\alpha-1}N^{-1}(1+o(1))$, and the identity
after~\eqref{eq:Phi} finishes the check. For $j$ bounded, $\Phi_b(j/N)\to1$ and (b)
reduces to (a). Without pure power tails the ratio $T_\minus(j)/T_\minus(N)$ keeps a slowly
varying factor in the range $1\ll j\ll N$, and (b) is the correct form.

\subsection{Proof of Theorem~\ref{thm:ends}(a)}

Let the start be fresh plus and fix $j$. We use Proposition~\ref{prop:final-run}. Fix
$r\le j$. We show that
$\E[T_\plus(N-j-S^\plus_r);S^\plus_r\le N-j-1]=T_\plus(N)(1+o(1))$. On
$\{S^\plus_r\le N/2\}$ the ratio $T_\plus(N-j-S^\plus_r)/T_\plus(N)$ is at most
$T_\plus(N/2-j)/T_\plus(N)\le C$ and tends to $1$ almost surely, so this part is
$T_\plus(N)(1+o(1))$ by dominated convergence. On $\{S^\plus_r>N/2\}$ put
$i=N-j-S^\plus_r$ and fix $I\ge1$. For fixed $r$ and $n\ge1$,
$u^\plus_r(n)\le r\sup_{s\ge n/r}f_\plus(s)$, because one of the $r$ summands is at least
$n/r$. Since $f_\plus(s)=T_\plus(s)-T_\plus(s+1)=o(T_\plus(s))$ by (R1), the part $i\le I$ is
at most $I\max_{n\ge N/2}u^\plus_r(n)=o(T_\plus(N))$. The part $i>I$ is at most
$T_\plus(I)\Prob(S^\plus_r>N/2)\le C_rT_\plus(I)T_\plus(N)$ by Lemma~\ref{lem:sums}(i) and
(R2). Let $N\to\infty$ and then $I\to\infty$. The second sum
in~\eqref{eq:final-explicit} has finitely many terms, each at most
$C_j\max_{r\le j}u^\plus_{r+1}(N-j)=o(T_\plus(N))$ by the same bound on $u^\plus_{r+1}$.
Hence $P_N(N-j)=T_\plus(N)\sum_{r\le j}u^\minus_r(j)+o(T_\plus(N))=T_\plus(N)u_\minus(j)+o(T_\plus(N))$.

For a fresh minus start we use Proposition~\ref{prop:first-run} with $\nu_\minus=f_\minus$.
Since $F^\plus_{N-d}(N-j)=0$ for $d>j$, it gives
\[
 P_N(N-j)=\sum_{1\le d\le j}f_\minus(d)\,F^\plus_{N-d}(N-j).
\] By the fresh plus case, $F^\plus_{N-d}(N-j)=T_\plus(N)u_\minus(j-d)+o(T_\plus(N))$
for each $d\le j$, and $\sum_{1\le d\le j}f_\minus(d)u_\minus(j-d)=u_\minus(j)$ for $j\ge1$
by the renewal equation. For $j=0$ the sum is empty. \qed

Only (RV) of $T_\plus$ is used in (a). The minus law enters only through $u_\minus$.

\subsection{Proof of Theorem~\ref{thm:ends}(b)}

Let the start be fresh plus. By (iii) of Lemma~\ref{lem:cd}, (LB$_\plus$) implies
(SRT$_\plus$), so Theorem~\ref{thm:local} applies. For $\varepsilon N\le j\le(1-\delta)N$
the claim follows from Theorem~\ref{thm:local} and~\eqref{eq:layer-bulk}, since
$u_\minus(j)\asymp\omega_\minus(j)$ there. For $j\le j_0$ with $j_0$ fixed it is (a), since
$\omega_\minus(j)\ge c_{j_0}$ and $|\Phi_b(j/N)-1|\to0$. Since $\Phi_b(y)=1+O(y^\alpha)$ and
$u_\minus(j)\ge c\,\omega_\minus(j)$ for $j\ge j_0$ under (SRT$_\minus$), it remains to
show the following. There is a function $\zeta$ with $\zeta(\eta)\to0$ as $\eta\to0$ such
that for every small $\eta>0$ there are $\varepsilon>0$ and $j_0$ with
\begin{equation}\label{eq:middle-layer}
 \limsup_{N\to\infty}\sup_{j_0\le j\le\varepsilon N}
 \Bigl|\frac{P_N(N-j)}{T_\plus(N)u_\minus(j)}-1\Bigr|\le\zeta(\eta),
\end{equation}
Let $k=N-j$ and write $P_N(k)=\mathrm I+\mathrm{II}$ as in Section~\ref{sec:local-proof}.
Put $M=\eta/T_\plus(N)$ and $\theta=\eta^{1/(2+2\alpha)}$. For small $\eta$ we have
$\eta^{1/\alpha}\le\theta/4$, since $1/\alpha>1>1/(2+2\alpha)$, so $a_M\le\theta N$ for large
$N$.

\emph{Term $\mathrm I$, $m\le M$.} We show that
$q^\plus_{m+1}(k)=T_\plus(N)(1+O(\eta^{1/2}+\theta+\varepsilon))$ uniformly in $m\le M$ and
$j\le\varepsilon N$. For the lower bound, the part $r<\theta N$ of~\eqref{eq:q-cover}
gives $q^\plus_{m+1}(k)\ge\Prob(S^\plus_m<\theta N)T_\plus(k)$, and
$\Prob(S^\plus_m\ge\theta N)\le CmT_\plus(\theta N)\le C\eta\theta^{-\alpha}\le C\eta^{1/2}$ by
Lemma~\ref{lem:sums}(i) and (R1). For the upper bound, the part $r<\theta N$ is at most
$T_\plus(k-\theta N)\le(1+o(1))(1-\varepsilon-\theta)^{-\alpha}T_\plus(N)$. For $r\ge\theta N$
we have $r\ge a_M\ge a_m$, so Lemma~\ref{lem:cd}(v) gives
$u^\plus_m(r)\le CmT_\plus(r)/r\le C\eta\theta^{-1-\alpha}/N$, and with (R3) this part is at
most $C\eta\theta^{-1-\alpha}N^{-1}\sum_{i\le N}T_\plus(i)\le C\eta\theta^{-1-\alpha}T_\plus(N)=C\eta^{1/2}T_\plus(N)$.

Next, the tail of the renewal sum. For $m\ge1/T_\minus(j)$ we have $a_m\ge cj$, so by
Lemma~\ref{lem:sums}(iii) for the minus sign and $T_\minus(j+1)\ge cT_\minus(j)$,
\[
 \sum_{m>M}u^\minus_m(j)\le\frac Cj\sum_{m>M}e^{-cmT_\minus(j)}
 \le\frac{C}{jT_\minus(j)}\,e^{-cMT_\minus(j)}
\]
whenever $MT_\minus(j)\ge1$. Now
$MT_\minus(j)\ge\eta T_\minus(\varepsilon N)/T_\plus(N)\to\eta b\varepsilon^{-\alpha}$. We choose
$\varepsilon=\varepsilon(\eta)\le\eta$ so small that $e^{-c\eta b\varepsilon^{-\alpha}/2}\le\eta$.
Since $u_\minus(j)\ge c/(jT_\minus(j))$ for $j\ge j_0$, this gives
$\sum_{m>M}u^\minus_m(j)\le C\eta\,u_\minus(j)$. Hence
$\sum_{1\le m\le M}u^\minus_m(j)q^\plus_{m+1}(k)=T_\plus(N)u_\minus(j)(1+O(\eta^{1/2}))$.

\emph{Term $\mathrm I$, $m>M$.} By Lemma~\ref{lem:key}, $q^\plus_{m+1}(k)\le CT_\plus(k)\le CT_\plus(N)$,
so this part is at most $CT_\plus(N)\sum_{m>M}u^\minus_m(j)\le C\eta\,T_\plus(N)u_\minus(j)$.

\emph{Term $\mathrm{II}$.} By Lemma~\ref{lem:decay}, $q^\minus_m(j)\le CT_\minus(j)e^{-cmT_\minus(j)}$.
Choose $c'>0$ with $a_{c'/T_\plus(N)}\le N/2$ for large $N$. For $m\le c'/T_\plus(N)$ we have
$k\ge N/2\ge a_m$, and Lemma~\ref{lem:cd}(v) gives $u^\plus_m(k)\le CmT_\plus(N)/N$. This part
is at most
\[
 C\,T_\minus(j)\frac{T_\plus(N)}N\sum_{m\ge1}me^{-cmT_\minus(j)}\le C\frac{T_\plus(N)}{NT_\minus(j)}
 \le C\frac jN\,T_\plus(N)u_\minus(j)\le C\varepsilon\,T_\plus(N)u_\minus(j).
\]
For $m>c'/T_\plus(N)$ we have $u^\plus_m(k)\le C/a_m\le C/N$, and this part is at most
$CN^{-1}e^{-c''s}$ with $s=T_\minus(j)/T_\plus(N)$, because $T_\minus(j)\sum_{m>c'/T_\plus(N)}e^{-cmT_\minus(j)}\le Ce^{-cc's}$.
Relative to $T_\plus(N)u_\minus(j)\ge c/(Nsj/N)$ this is at most
$C(j/N)\,s\,e^{-c''s}\le C\varepsilon$.

Collecting the terms gives~\eqref{eq:middle-layer} with
$\zeta(\eta)=C(\eta^{1/2}+\theta+\varepsilon+\eta)\le C\eta^{1/(2+2\alpha)}$.
For a fresh minus start, Proposition~\ref{prop:fresh-identity} gives
$P_N(k)=\sum_{m\ge1}[u^\plus_m(k)q^\minus_{m+1}(j)+u^\minus_m(j)q^\plus_m(k)]$. The second sum
is treated as $\mathrm I$ and gives $T_\plus(N)\sum_{m\ge1}u^\minus_m(j)=T_\plus(N)u_\minus(j)$ for
$j\ge1$, and the first is treated as $\mathrm{II}$. \qed

\subsection{Proof of Theorem~\ref{thm:ends}(c)}

By (iii) of Lemma~\ref{lem:cd}, (LB$_\pm$) implies (SRT$_\pm$). Apply (b) at the plus end
with $\delta=1/2$. For $J_N\le j\le N/2$, (SRT$_\minus$) gives $u_\minus(j)\asymp\omega_\minus(j)$,
so the error $o(\omega_\minus(j))$ is $o(u_\minus(j))$. With pure power tails,
$T_\plus(N)u_\minus(j)=\frac{\sin\pi\alpha}\pi\frac{c_+N^{-\alpha}}{c_-j^{1-\alpha}}(1+o(1))
=\frac{\sin\pi\alpha}{\pi b}\,N^{-1}(j/N)^{\alpha-1}(1+o(1))$ uniformly in $j\ge J_N$, with
$b=c_-/c_+$. Multiplying by $\Phi_b(j/N)$ gives $N^{-1}\La(1-j/N)(1+o(1))$. The minus end is
the same with the signs exchanged. \qed

\subsection{An upper bound under the strong renewal theorem}\label{sec:lemE}

Part (b) of Theorem~\ref{thm:ends} uses (LB$_\plus$). Its upper half, the bound
$P_N(N-j)\le CT_\plus(N)\omega_\minus(j)$ on $j\le N/2$, holds under the strong renewal theorem
alone, except on the line $\alpha=1/2$, where we need the slowly varying factor to be
bounded. Section~\ref{sec:aging} uses this bound for the remainder walk after a long initial
run.

\begin{lemma}[the end layer from above]\label{lem:E}
Assume (BAL) and (SPAN), and one of the following: $\alpha>1/2$; $\alpha<1/2$ and (SRT$_\pm$);
$\alpha=1/2$, (SRT$_\pm$), and $nT_\plus(n)^2$ bounded above and away from zero. Then there is
$C$ such that for all $n\ge1$ and $0\le i\le n/2$,
\[
 F^\minus_n(i)\le C\,T_\plus(n)\,\omega_\plus(i),\qquad F^\plus_n(n-i)\le C\,T_\plus(n)\,\omega_\minus(i).
\]
\end{lemma}

\begin{proof}
We prove the first bound, and the second is its mirror. At $i=0$, $F^\minus_n(0)=T_\minus(n)\le CT_\plus(n)$.
Let $i\ge1$ and $j=n-i\ge n/2$. By Proposition~\ref{prop:fresh-identity} for a fresh minus start,
$F^\minus_n(i)=\sum_mu^\plus_m(i)q^\minus_{m+1}(j)+\sum_mu^\minus_m(j)q^\plus_m(i)$. In the first sum
Lemma~\ref{lem:key} gives $q^\minus_{m+1}(j)\le CT_\minus(j)\le CT_\plus(n)$, and the sum of
$u^\plus_m(i)$ over $m$ is $u_\plus(i)\le C\omega_\plus(i)$ by (SRT$_\plus$), which is automatic for
$\alpha>1/2$. In the second sum Lemma~\ref{lem:decay} gives
$q^\plus_m(i)\le CT_\plus(i)e^{-cmT_\plus(i)}$. If $\alpha>1/2$, Lemma~\ref{lem:cd}(ii) gives
$u^\minus_m(j)\le CmT_\minus(j)/a_m\le CT_\plus(n)m/a_m$, and since $m/a_m$ is regularly varying with
index $1-1/\alpha\in(-1,0)$, $\sum_m(m/a_m)e^{-cmT_\plus(i)}\le CM^2/a_M$ with $M=1/T_\plus(i)$, and
$a_M\asymp i$, so the second sum is at most $CT_\plus(n)/(iT_\plus(i))$. If $\alpha<1/2$, we bound
$e^{-cmT_\plus(i)}$ by $1$, so the second sum is at most
$CT_\plus(i)u_\minus(j)\le CT_\plus(i)/(nT_\plus(n))$ by (SRT$_\minus$), and
$T_\plus(i)/(nT_\plus(n))\le CT_\plus(n)/(iT_\plus(i))$ because $iT_\plus(i)^2\le CnT_\plus(n)^2$ for
$i\le n$, as $xT_\plus(x)^2$ is regularly varying with the positive index
$1-2\alpha$~\cite[Theorem~1.5.3]{bgt1989}. If $\alpha=1/2$, the same route works when $xT_\plus(x)^2$ is
bounded above and away from zero.
\end{proof}

\subsection{The global mode}

\begin{corollary}[the global mode]\label{cor:mode}
Assume (BAL), (SPAN) and (LB$_\pm$), and let the start be fresh of either sign.
Let $\kappa_\plus(j)=\lim_NP_N(N-j)/T_\plus(N)$ and $\kappa_\minus(j)=\lim_NP_N(j)/T_\minus(N)$,
given by Theorem~\ref{thm:ends}(a). There is $J$ such that for all large $N$ every global
mode of $P_N$ lies in $\{0,\dots,J\}\cup\{N-J,\dots,N\}$. If moreover
$\max_j\kappa_\plus(j)>b\max_j\kappa_\minus(j)$ and the maximum of $\kappa_\plus$ is attained at a
single $j_+$, then the global mode is $N-j_+$ for all large $N$.
\end{corollary}

\begin{proof}
Both profiles tend to $0$ and are not identically $0$, so their maxima are positive and
attained. By Theorem~\ref{thm:local}, $P_N\le C/N$ on $[\varepsilon N,(1-\varepsilon)N]$. By
Theorem~\ref{thm:ends}(b), $P_N(N-j)\le CT_\plus(N)\omega_\minus(j)\le CT_\plus(N)\omega_\minus(J)$
for $J\le j\le\varepsilon N$, and the same at the minus end. Since $NT_\plus(N)\to\infty$ and
$\omega_\pm(J)\to0$ as $J\to\infty$, for $J$ large and then $N$ large every value of $P_N$
on $[J,N-J]$ is less than half of
$\min\{T_\plus(N)\max\kappa_\plus,\,T_\minus(N)\max\kappa_\minus\}$. The values at distance
less than $J$ from the ends converge after division by $T_\plus(N)$ or $T_\minus(N)$, and
$T_\minus(N)\sim bT_\plus(N)$. This gives both claims.
\end{proof}

For a fresh plus start, $\kappa_\plus=u_\minus$ has $\kappa_\plus(0)=1$, and the mode is at
an end for every $\alpha\in(0,1)$. For $1<\alpha<2$ the comparison between the atoms and the
bulk has a golden ratio threshold~\cite{paperF}. Here the end always wins, because the
bulk values are of order $1/N$ and the end values of order $T(N)\gg1/N$.

\begin{remark}[what (b) needs]\label{rem:M0}
We use (LB$_\plus$) only through Lemma~\ref{lem:cd}(v), in the range $j\le\varepsilon N$. Under
(SRT$_\pm$) alone Lemma~\ref{lem:E} gives the upper half of (b) when $\alpha\ne1/2$, and we do
not know whether the lower half holds. It would follow from the bound
\[
 \sum_{i\le\theta N}T_\plus(i)\sum_{m\le\eta/T_\plus(N)}u^\minus_m(j)\,u^\plus_m(N-j-i)\le\zeta(\theta,\eta)\,T_\plus(N)u_\minus(j)
\]
for $j\le\varepsilon N$, with $\zeta(\theta,\eta)\to0$ as $\theta,\eta\to0$. The bound of
Lemma~\ref{lem:cd}(ii) loses a factor of order $N/j$ here.
\end{remark}
